\section{Resultados e Discussão} \label{section: results}

\subsection{Validação do Protocolo por \textit{Redocking}}
PREENCHER: RMSD do ligante cocristalizado (meta $\le 2,0\text{ \AA}$) e, se pertinente, o efeito da protonação do controle. Figura sugerida: sobreposição cristal $\times$ pose.

\subsection{Escores e Eficiência do Ligante}
PREENCHER: tabela com todas as métricas lado a lado ($\Delta G_{\text{Vina}}$, $\text{CNN}_{\text{affinity}}$, $\text{CNN}_{\text{score}}$, $\text{LE}_{\text{Vina}}$, $d_{\text{sítio}}$, $N_{\text{válidas}}/N$). Indique a métrica pré-registrada, se a posição do composto líder muda conforme a métrica e se a amplitude dos escores é suficiente para discriminar a série.

\subsection{Modo de Ligação e Interações}
PREENCHER: posição de cada ligante na cavidade e interações tipificadas (resíduo, tipo, distância, ângulo e fragmento do ligante). Figuras sugeridas: diagrama 2D de interações e figura 3D do sítio.

\subsection{Relação Estrutura--Atividade}
PREENCHER: o que diferencia os compostos ativos dos menos ativos segundo o modelo, sempre como hipótese estrutural.
